\hypertarget{project-kennel-construction-principles}{%
\chapter{Project Kennel: Construction
Principles}\label{project-kennel-construction-principles}}

These are the postures the code and architecture hold to, sorted in
three small categories: Stance, Supply and Construction. Where the
design register states what the system is and why, these state how its
code is built, under a reader who does not trust the author. They are
construction, not design.

Each rests on a section of the coding standards, cited as §N. The
standards state the rules; this register states the postures they serve.
Where a construction principle is a design principle realised in code,
the design slug is named: the two registers meet at the altitude line,
posture above and mechanism below.

What is not here is the working discipline: how a claim carries its
evidence, how a change earns its place, how the tests are written before
the implementation they check, how a contradiction between the substrate
and the corpus is surfaced rather than patched. Those govern the people
and the process, not the landed artefact, and they live in the coding
standards and the standing orders. This register holds only what shows
up in the code and the architecture.

\hypertarget{stance}{%
\section{Stance}\label{stance}}

\hypertarget{read-by-the-hostile}{%
\subsection{read-by-the-hostile}\label{read-by-the-hostile}}

The reader does not trust the author. Code in the repository is read
line by line by people who assume it is hostile until they have
satisfied themselves it is not, and the bar it is held to is OpenSSH and
libpam: auditable in full, by a stranger, with the author absent to
explain it (§1). This is \texttt{good-boy} turned on the code itself. A
design that refuses to ask the workload to be trustworthy cannot be
implemented by code that asks the same of its reader. So the code is
written to be checked, not believed: small enough to read, plain enough
to follow, arranged so a reviewer can confirm a property rather than
take it on faith. Cleverness that buys a line and costs an hour of
reading is a net loss, because the hour is paid by everyone who audits
the system and the line is saved once.

\hypertarget{built-to-outlast}{%
\subsection{built-to-outlast}\label{built-to-outlast}}

The code must compile, test, and run the same way on the next
long-term-support distribution five years from now (§2). The toolchains
are pinned, the build is offline against vendored sources, and release
builds are reproducible hash for hash on two independent runners, so
that what shipped can be rebuilt and checked long after the machine that
built it is gone. A confinement tool an auditor cannot reproduce is one
an auditor cannot trust, and trust that cannot be re-established on a
fresh machine is trust with an expiry date. The cost of writing against
a moving substrate is paid once by the author; the cost of a build that
cannot be reproduced is paid by every future reader, so the trade runs
one way.

\hypertarget{supply}{%
\section{Supply}\label{supply}}

\hypertarget{own-your-supply-chain}{%
\subsection{own-your-supply-chain}\label{own-your-supply-chain}}

Every dependency is attack surface and audit burden, and the default
answer to taking one is no (§5). What the standard library and the
existing tree already provide is tried first; the question is not
whether a crate exists but whether the thing could be written in a day
and maintained without embarrassment. What is taken is pinned to an
exact version, vendored as the audited bytes, and checked against an
independently computed manifest rather than a hash the registry vouches
for, because trust on first use against a registry no one controls is
not a defence. This is \texttt{do-less} on the supply chain: the smaller
the set of things outside the tree the system must trust, the smaller
the set that can betray it. A hand-rolled loader that replaces a
thousand vendored files is not reinvention; it is the refusal of a
dependency whose cost outweighs its saving.

\hypertarget{dont-roll-your-own}{%
\subsection{dont-roll-your-own}\label{dont-roll-your-own}}

The opposite error is hubris, and the discipline refuses it as hard.
Cryptography, transport security, compression, character-set decoding,
the hardened parsers where a subtle bug is an exploit: these are
minefields cleared over decades by specialists, and a fresh
implementation forfeits the thousands of hours and eyes the existing one
already carries (§5.1). Here the tree admits the dependency. The line
between this and \texttt{own-your-supply-chain} is one question asked
plainly: can we audit our own version in full, and would it earn as many
eyes as the thing it would replace. Yes to both and the dependency is
bloat; no to either and reimplementing is recklessness. The two halves
are one rule reading its inputs, which is why a small reviewable loader
and a refusal to write our own signature scheme are not a contradiction.

The project's one crypto primitive shows the pair deciding together. A
solid security primitive was needed and was never going to be written
here, so the only question was which to depend on, and the answer was
the smallest vetted one. Among the candidates, rustls was about 35,000
lines on its own and hundreds of thousands once its dependencies were
counted; a sha256 crate was under 400 lines in itself but close to 8,000
with its tree; ed25519 came to about 800 lines all in, and won on total
footprint, the least code to read that still does the job. The gap
between each pair of figures is the lesson: the size of the thing you
came for tells you nothing, because a dependency costs the whole tree it
drags in, not the leaf. The direction is worth marking too: ed25519 is
in the tree because supply chose the smallest vetted primitive, and the
trust material the design register spends is spending what supply
admitted, rather than the reverse.

\hypertarget{construction}{%
\section{Construction}\label{construction}}

\hypertarget{parse-dont-validate}{%
\subsection{parse-dont-validate}\label{parse-dont-validate}}

Every byte that enters from outside is untrusted until parsing has
produced a typed value, and the parse is the validation, not a check
bolted on after (§10). A typed value is evidence that the checks the
type requires have run; raw input has no standing until it has become
something the rest of the program is allowed to hold. Configuration is
not exempt and is the sharp case, because a file looks like an internal
artefact and draws less scrutiny than network bytes while arriving from
templates fetched over the wire, deltas hand-written, caches that may
have been tampered with, and the very agents the system exists to
confine. Unknown fields are rejected, reads are bounded before they
begin, and the crossing is visible in the code as a parser call rather
than a silent conversion. This is \texttt{mediate-use-not-reach} and
\texttt{construction-by-absence} at the type boundary: a constructed
value is a mediated one, and what was never parsed was never admitted.

\hypertarget{make-invalid-unrepresentable}{%
\subsection{make-invalid-unrepresentable}\label{make-invalid-unrepresentable}}

Where a value carries security meaning it is a newtype whose constructor
does the validation, and where an invalid state is forbidden by
invariant it is removed from the type rather than checked at runtime
(§11). A forbidden flag is not a boolean someone must remember to test;
it is a shape that cannot be built. The constructed type is the proof
the check happened, so the proof cannot be skipped without the value
failing to exist. This is \texttt{construction-by-absence} carried into
the type system: the same move that makes a withheld resource absent
from a workload's world makes an invalid state absent from the
program's, and in both the gain is that there is nothing left to guard,
because there is nothing left to be wrong.

\hypertarget{errors-not-panics}{%
\subsection{errors-not-panics}\label{errors-not-panics}}

Shipped code has no unhandled path (§8). There is no \texttt{unwrap}, no
\texttt{panic!}, no \texttt{todo!} in anything that ships; every
fallible thing returns a typed error specific enough to act on; and the
one permitted \texttt{expect} carries, in its message, the invariant
that makes it infallible, so the claim of safety is written where a
reader can check it. At a trust boundary an error is an expected
outcome, not an exceptional one, given a named variant and an audit
record rather than a crash. This is \texttt{refuse-to-start} in the
small: a component that meets a condition it cannot honour stops cleanly
and says so, rather than proceeding on a value it had to assume. The
privhelper goes further and aborts rather than unwind, because unwinding
through cleanup it cannot trust is a worse failure than stopping.

\hypertarget{quarantine-the-unsafe}{%
\subsection{quarantine-the-unsafe}\label{quarantine-the-unsafe}}

\texttt{unsafe} and C are quarantined into a few small, single-purpose
crates, each sized to be read in a single sitting, and every other crate
forbids \texttt{unsafe} outright (§3, §4). The syscall and BPF glue that
no safe abstraction covers is concentrated where two maintainers must
read every block, not spread thin across the tree where each instance
escapes notice. This is the design register's \texttt{rule-of-1}
realised in the workspace: the hazard that cannot be removed is gathered
into the few crates that carry it and nothing else, each holding one
narrow purpose, so a bug in the unsafe code is bounded by the small
crate that holds it. It composes with \texttt{dont-roll-your-own}: lean
on a vetted crate to avoid the unsafe you do not have to write, and
quarantine the irreducible remainder no crate covers. Avoid what you
can, gather what you cannot.
